\sect{भाद्रपद-शुक्ल-पद्मा-एकादशी-माहात्म्यम्}
\label{sec:padma-bhadrapada-shukla-padma}

\ganapatyadiDhyanam
\padmaPuranam

\dnsub{अथ एकोनषष्टितमोऽध्यायः॥५९}

\uvacha{युधिष्ठिर उवाच}

\twolineshloka
{नभस्यस्य सिते पक्षे किन्नामैकादशी भवेत्}
{को देवः को विधिस्तस्य एतदाख्याहि केशव}% ॥१॥

\uvacha{श्रीकृष्ण उवाच}

\twolineshloka
{कथयामि महीपाल कथामाश्चर्यकारिणीम्}
{कथयामास यां ब्रह्मा नारदाय महात्मने}% ॥२॥

\uvacha{नारद उवाच}

\threelineshloka
{कथयस्व प्रसादेन चतुर्मुख नमोऽस्तु ते}
{नभस्य शुक्लपक्षे तु किं नामैकादशी भवेत्}
{एतदिच्छाम्यहं श्रोतुं विष्णोराराधनाय वै}% ॥३॥

\uvacha{ब्रह्मोवाच}

\twolineshloka
{वैष्णवोऽसि मुनिश्रेष्ठ साधुपृष्टं किल त्वया}
{नातः परतरा लोके पवित्रा हरिवासरात्}% ॥४॥

\twolineshloka
{पद्मा नामेति विख्याता नभस्यैकादशी सिता}
{हृषीकेशः पूज्यतेऽस्यां कर्तव्यं व्रतमुत्तमम्}% ॥५॥

\twolineshloka
{कथयामि तवाग्रेऽहं कथां पौराणिकीं शुभाम्}
{यस्याः श्रवणमात्रेण महापापं प्रणश्यति}% ॥६॥

\twolineshloka
{मान्धाता नाम राजर्षिर्विवस्वद्वंशसम्भवः}
{बभूव चक्रवर्ती स सत्यसन्धः प्रतापवान्}% ॥७॥

\twolineshloka
{धर्मतः पालयामास प्रजाः पुत्रानिवौरसान्}
{न तस्य राज्ये दुर्भिक्षं नाधयो व्याधयस्तथा}% ॥८॥

\twolineshloka
{निरातङ्काः प्रजास्तस्य धनधान्यसमेधिताः}
{न्यायेनोपार्जितं वित्तं तस्य कोशे महीपते}% ॥९॥

\twolineshloka
{स्वस्वधर्मे प्रवर्तन्ते सर्वे वर्णाश्रमास्तथा}
{कामधेनुसमाभूमिस्तस्य राज्ये महीपतेः}% ॥१०॥

\twolineshloka
{तस्यैवं कुर्वतो राज्यं बहुवर्षगणा गताः}
{अथैकस्मिंश्च सम्प्राप्ते विपाकः कर्मणः खलु}% ॥११॥

\twolineshloka
{वर्षत्रयं तद् विषये न ववर्ष बलाहकः}
{तेन भग्नाः प्रजास्तस्य बभूवुः क्षुधयार्दिताः}% ॥१२॥

\threelineshloka
{स्वाहा स्वधा वषट्कार वेदाध्ययनवर्जिताः}
{बभूव विषयस्तस्या भाग्येन दैवपीडितः}
{अथ प्रजाः समागम्य राजानमिदमब्रुवन्}% ॥१३॥

\uvacha{प्रजा ऊचुः}

\twolineshloka
{श्रोतव्यं नृपशार्दूल प्रजानां वचनं त्वया}
{आपो नारा इति प्रोक्ताः पुराणेषु मनीषिभिः}% ॥१४॥

\twolineshloka
{अयनं भगवतस्तस्मान्नारायण इति स्मृतः}
{पर्जन्यरूपो भगवान् विष्णुः सर्वगतः स्थितः}% ॥१५॥

\threelineshloka
{स एवं कुरुते वृष्टिं वृष्टेरन्नं ततः प्रजाः}
{तदभावे नृपश्रेष्ठ क्षयं गच्छन्ति वै प्रजाः}
{तथा कुरु नृपश्रेष्ठ योगः क्षेमो यथा भवेत्}% ॥१६॥

\uvacha{राजोवाच}

\twolineshloka
{सत्यमुक्तं भवद्भिश्च न मिथ्याभिहितं क्वचित्}
{अन्नं ब्रह्म यतः प्रोक्तमन्ने सर्वं प्रतिष्ठितम्}% ॥१७॥

\twolineshloka
{अन्नाद् भवन्ति भूतानि जगदन्नेन वर्तते}
{इत्येवं श्रूयते लोके पुराणे बहुविस्तरे}% ॥१८॥

\twolineshloka
{नृपाणामपचारेण प्रजानां पीडनं भवेत्}
{नाहं पश्याम्यात्मकृतमेवं बुद्ध्या विचारयन्}% ॥१९॥

\twolineshloka
{तथाऽपि प्रयतिष्यामि प्रजानां हितकाम्यया}
{इति कृत्वा मतिं राजा परिमेयपरिच्छदः}% ॥२०॥

\twolineshloka
{नमस्कृत्य विधातारं जगाम गहनं वनम्}
{चचार मुनिमुख्यांश्च आश्रमान्तापसैः श्रितान्}% ॥२१॥

\twolineshloka
{ददर्शाथ ब्रह्मसुतमृषिमाङ्गिरसं नृपः}
{तेजसा द्योतितदिशं द्वितीयमिव पद्मजम्}% ॥२२॥

\twolineshloka
{तं दृष्ट्वा हर्षितो राजा अवतीर्य स्ववाहनात्}
{नमश्चक्रेऽस्य चरणौ कृताञ्जलिपुटो वशी}% ॥२३॥

\twolineshloka
{मुनिस्तमभिनन्द्याथ स्वस्तिवाचनपूर्वकम्}
{पप्रच्छ कुशलं राज्ये सप्तस्वङ्गेषु भूपतेः}% ॥२४॥

\twolineshloka
{निवेदयित्वा कुशलं पप्रच्छानामयं नृपः}
{दत्तासनो गृहीतार्घ्य उपविष्टोऽस्य सन्निधौ}% ॥२५॥

\onelineshloka
{प्रत्युवाच मुनिं राजा पृष्टो ह्यागमकारणम्}% ॥२६॥

\uvacha{राजोवाच}

\twolineshloka
{भगवन् धर्मविधिना मम पालयतो महीम्}
{अनावृष्टिश्च संवृत्ता नाहं वेद्म्यत्र कारणम्}% ॥२७॥

\twolineshloka
{संशयच्छेदनायात्र आगतोऽहं तवान्तिके}
{योगक्षेमविधानेन प्रजानां कुरु निर्वृतिम्}% ॥२८॥

\uvacha{ऋषिरुवाच}

\twolineshloka
{एतत्कृतयुगं राजन् युगानामुत्तमं मतम्}
{अत्र ब्रह्मपरा लोका धर्मश्चात्र चतुष्पदः}% ॥२९॥

\twolineshloka
{अस्मिन् युगे तपोयुक्ता ब्राह्मणा नेतरा जनाः}
{विषये तव राजेन्द्र वृषलोऽयं तपस्यति}% ॥३०॥

\twolineshloka
{एतस्मात् कारणाच्चैव न वर्षति बलाहकः}
{कुरु तस्य वधे यत्नं येन दोषः प्रशाम्यति}% ॥३१॥

\uvacha{राजोवाच}

\twolineshloka
{नाहमेनं वधिष्यामि तपस्यन्तमनागसम्}
{धर्मोपदेशं कथय उपसर्गविनाशनम्}% ॥३२॥

\uvacha{ऋषिरुवाच}

\twolineshloka
{यद्येवं तर्हि नृपते कुरुष्वैकादशीव्रतम्}
{नभस्यस्य सिते पक्षे पद्मा नामेति विश्रुता}% ॥३३॥

\twolineshloka
{तस्या व्रतप्रभावेन सुवृष्टिर्भविता ध्रुवम्}
{सर्वसिद्धिप्रदा ह्येषा सर्वोपद्रवनाशिनी}% ॥३४॥

\twolineshloka
{अस्या व्रतं कुरु नृप सप्रजः सपरिच्छदः}
{इति वाक्यमृषेः श्रुत्वा राजा स्वगृहमागतः}% ॥३५॥

\twolineshloka
{भाद्रमासे सिते पक्षे पद्माव्रतमथाकरोत्}
{प्रजाभिः सह सर्वाभिश्चातुर्वर्ण्यसमन्वितः}% ॥३६॥

\twolineshloka
{एवं व्रते कृते राजन् प्रववर्ष बलाहकः}
{जलेन प्लाविता भूमिरभवत् सस्यशालिनी}% ॥३७॥

\twolineshloka
{ऋषीश्वरप्रभावेन लोकाः सौख्यं प्रपेदिरे}
{एतस्मात् कारणादेवं कर्तव्यं व्रतमुत्तमम्}% ॥३८॥

\twolineshloka
{दध्योदनयुतं तस्यां जलपूर्णं घटं द्विजे}
{वस्त्रसंवेष्टितं दत्त्वा छत्रोपानहमेव च}% ॥३९॥

\twolineshloka
{नमो नमस्ते गोविन्द बुधश्रवणसंज्ञक}
{अघौघसङ्क्षयं कृत्वा सर्वसौख्यप्रदो भव}% ॥४०॥

\twolineshloka
{भुक्तिमुक्तिप्रदश्चैव लोकानां सुखदायकः}
{पठनाच्छ्रवणाद् राजन् सर्वपापैः प्रमुच्यते}% ॥४१॥

॥इति श्रीपाद्मे महापुराणे पञ्चपञ्चाशत्साहस्र्यां संहितायामुत्तरखण्डे उमापतिनारदसंवादान्तर्गतकृष्णयुधिष्ठिरसंवादे भाद्रपद-शुक्ल-पद्मा-एकादशी-माहात्म्यं नाम एकोनषष्टितमोऽध्यायः॥५९॥

\genMangalaShloka

\hyperref[sec:ekadashi_mahatmyam_padma_puranam]{\closesub}
\clearpage

